% Guía del Taller virtual ATA EPE: cómo se estudia, se construye y se publica cada módulo.
% Compilar:  latexmk -pdf guia_del_taller.tex
\documentclass[11pt,a4paper]{article}
% Preámbulo común de las fichas de estudio del Taller virtual ATA EPE.
\usepackage[spanish,es-noshorthands,es-tabla]{babel}
\usepackage[utf8]{inputenc}
\usepackage[T1]{fontenc}
\usepackage{lmodern}
\usepackage[margin=2.2cm]{geometry}
\usepackage[table]{xcolor}
\usepackage{amsmath,amssymb,bm}
\usepackage{tabularx,array,booktabs}
\usepackage{enumitem}
\usepackage{fontawesome5}
\usepackage[most]{tcolorbox}
\usepackage[hidelinks]{hyperref}
\usepackage{listings}

\definecolor{morado}{HTML}{2A1B7E}
\definecolor{amarillo}{HTML}{FFCF2E}
\definecolor{rojo}{HTML}{C4142C}
\definecolor{verde}{HTML}{1F9D55}
\definecolor{gris}{HTML}{5D5A73}

\newcolumntype{Y}{>{\raggedright\arraybackslash}X}
\renewcommand{\arraystretch}{1.25}
\setlist{itemsep=1pt,topsep=3pt}
\numberwithin{equation}{section}
\lstset{basicstyle=\ttfamily\small,columns=fullflexible,keepspaces=true,frame=single,rulecolor=\color{gris!40},
  literate={á}{{\'a}}1 {é}{{\'e}}1 {í}{{\'i}}1 {ó}{{\'o}}1 {ú}{{\'u}}1 {ñ}{{\~n}}1 {Ω}{{$\Omega$}}1
    {θ}{{$\theta$}}1 {λ}{{$\lambda$}}1 {Δ}{{$\Delta$}}1 {μ}{{$\mu$}}1 {·}{{$\cdot$}}1 {←}{{$\leftarrow$}}1 {§}{{\S}}1 {→}{{$\to$}}1 {≠}{{$\neq$}}1 {É}{{\'E}}1 {Á}{{\'A}}1 {∨}{{$\vee$}}1 {∧}{{$\wedge$}}1 {σ}{{$\sigma$}}1 {×}{{$\times$}}1 {≈}{{$\approx$}}1}

\newcommand{\rev}[1]{\widetilde{#1}}
\newcommand{\I}{\mathbf{I}}
\newcommand{\vect}[1]{\bm{#1}}
\newcommand{\A}{\mathcal{A}}

\newtcolorbox{meta}{colback=morado,colframe=morado,coltext=white,arc=3mm}
\newtcolorbox{lectura}{colback=morado!5,colframe=morado,fonttitle=\bfseries,title={\faBook~Lecturas}}
\newtcolorbox{ejercicio}[1]{colback=amarillo!10,colframe=amarillo!80!black,fonttitle=\bfseries,title={\faPencil*~Ejercicio #1}}
\newtcolorbox{control}{colback=rojo!5,colframe=rojo,fonttitle=\bfseries,title={\faQuestionCircle~Preguntas de control}}
\newtcolorbox{codigo}[1][De la ecuación al código]{colback=verde!6,colframe=verde,fonttitle=\bfseries,title={\faCode~#1}}
\newtcolorbox{cierre}{colback=verde!6,colframe=verde,fonttitle=\bfseries,title={\faFlagCheckered~Cierre del módulo}}
\newtcolorbox{nota}[1]{colback=amarillo!12,colframe=amarillo!80!black,fonttitle=\bfseries,title={\faLightbulb~#1}}

% \cabecera{número}{título}{subtítulo}
\newcommand{\cabecera}[3]{%
  \title{\textcolor{morado}{\textbf{Ficha #1. #2}}\\[3pt]\large Taller virtual ATA EPE · #3}
  \author{}\date{}}

\title{\textcolor{morado}{\textbf{Guía del taller}}\\[3pt]
  \large Taller virtual ATA EPE · Cómo se estudia, se construye y se publica}
\author{}\date{}

\begin{document}
\maketitle
\vspace{-2em}

\begin{meta}
\textbf{\textcolor{amarillo}{Para qué es esta guía.}} El taller es un curso que se aprende construyendo: nueve
módulos, cada uno con una ficha de estudio, unas pruebas y una función nueva en la web. Esta guía explica el ciclo de
trabajo de cada módulo, las herramientas, las reglas del código y cómo se sube y se publica lo hecho.
\end{meta}

\tableofcontents

% =================================================================================================
\section{El mapa del curso}

\begin{tabularx}{\linewidth}{@{}l l Y l@{}}
\toprule
& \textbf{Ficha} & \textbf{Tema} & \textbf{Horas}\\
\midrule
0 & \texttt{ficha00\_pga} & Del producto geométrico a la PGA: el lenguaje del simulador & 15\\
1 & \texttt{ficha01\_cinematica} & Movimientos rígidos, tornillos, velocidad, armar e imantar piezas & 15\\
2 & \texttt{ficha02\_dinamica} & Masa, inercia, Newton–Euler en PGA, integradores en el grupo & 20\\
3 & \texttt{ficha03\_analitica\_i} & Ligaduras, grados de libertad, ejes, rango del jacobiano & 15\\
4 & \texttt{ficha04\_ligaduras} & Uniones dinámicas, tornillos y tuercas, cuerpos compuestos & 25\\
5 & \texttt{ficha05\_contacto} & Colisiones, contacto unilateral y fricción de Coulomb & 35\\
6 & \texttt{ficha06\_transmisiones} & Ruedas de fricción, bandas elásticas, manivela & 20\\
7 & \texttt{ficha07\_analitica\_ii} & Lagrangiano reducido, trabajo virtual, «Ver las cuentas» & 15\\
8 & \texttt{ficha08\_compartir} & Guardar, compartir, retos y realidad aumentada & 10\\
\midrule
& & \textbf{Total} & \textbf{$\approx$ 170 h}\\
\bottomrule
\end{tabularx}

\medskip
Las horas son una estimación para alguien que estudia y programa a la vez. A un ritmo de 6 a 8 horas por semana, el
curso completo toma de 5 a 7 meses. El orden importa: cada módulo usa lo de los anteriores.

% =================================================================================================
\section{El ciclo de cada módulo}

\begin{tabularx}{\linewidth}{@{}l l Y@{}}
\toprule
\textbf{Paso} & & \textbf{Qué se hace}\\
\midrule
1 & \faBook~Leer & Las lecturas de la ficha, en el orden indicado. Lo marcado en negrita es lo esencial.\\
2 & \faPencil*~Hacer & Los ejercicios a mano, en papel o en \texttt{estudio/soluciones/fichaNN.tex}.\\
3 & \faQuestionCircle~Responder & Las preguntas de control, por escrito. Si una no sale, se vuelve a las lecturas.\\
4 & \faProjectDiagram~Discretizar & Escribir en pseudocódigo el algoritmo del módulo, con sus fuentes de error.\\
5 & \faFileContract~Contrato & Se generan las pruebas y el esqueleto del módulo (con Claude, a partir de la
  ficha y de lo ya implementado).\\
6 & \faCode~Implementar & En una rama propia, hasta que pasen las pruebas.\\
7 & \faCheck~Verificar & Pruebas automáticas, prueba visual en el navegador y revisión del código.\\
8 & \faUpload~Subir & Fusionar la rama con \texttt{main} y subir: GitHub Pages publica solo.\\
9 & \faFlagCheckered~Cerrar & Marcar el módulo en la bitácora y anotar lo aprendido y lo pendiente.\\
\bottomrule
\end{tabularx}

\begin{nota}{La regla de oro}
No se escribe código de un módulo sin haber hecho los pasos 1 a 4. Si al programar aparece algo que no se entiende,
se vuelve a la ficha: el código es la última expresión de una idea que ya está clara en papel.
\end{nota}

% =================================================================================================
\section{Herramientas}

\begin{tabularx}{\linewidth}{@{}l Y@{}}
\toprule
\textbf{Herramienta} & \textbf{Para qué}\\
\midrule
Node 22 & Correr las pruebas: \texttt{node --test "pruebas/*.test.js"} desde \texttt{taller/}.\\
Navegador & Ver la cara pública. Servidor local desde la carpeta \texttt{Cartilla}:
  \texttt{python3 -m http.server 8765}, y abrir \texttt{http://localhost:8765/taller/}.\\
Git y GitHub & Ramas por módulo, historial y publicación (el repositorio de la web).\\
LaTeX & Las fichas y las soluciones: \texttt{latexmk -pdf fichaNN\_*.tex}.\\
Editor & Cualquiera con resaltado de JavaScript (VS Code, por ejemplo).\\
\bottomrule
\end{tabularx}

% =================================================================================================
\section{Carpetas}

\begin{lstlisting}
taller/
  plan/          el plan general (modelo físico, módulos, bibliografía)
  estudio/       esta guía, las fichas, comun.tex, soluciones/ y bitacora.md
  nucleo/        el motor: pga, cinemática, cuerpo, integrador, ligaduras, colisión, contacto
  kit/           lo propio del kit: catálogo con puertos, ensamble, tornillería, ruedas, bandas
  analisis/      grados de libertad, lagrangiano reducido, trabajo virtual
  app/           la cara pública: escena, interfaz, instrumentos
  pruebas/       una prueba por archivo del núcleo, del kit y del análisis
  index.html     la entrada pública del taller
\end{lstlisting}

% =================================================================================================
\section{Reglas del código}

\begin{enumerate}[leftmargin=*]
  \item \textbf{Nombres en español}, como el resto de la web: \texttt{velocidad}, \texttt{motor}, \texttt{inercia}.
  \item \textbf{Cada función dice qué ecuación implementa}: «ficha 2, ec. 3.4». Si no hay ecuación, falta teoría.
  \item \textbf{El núcleo no conoce el dibujo.} \texttt{nucleo/}, \texttt{kit/} y \texttt{analisis/} no importan
        three.js ni tocan la página; así se prueban con Node. Solo \texttt{app/} dibuja.
  \item \textbf{Sin librerías de física ni de álgebra.} Todo lo numérico se escribe aquí.
  \item \textbf{Multivectores en \texttt{Float64Array(16)}}, como en la ficha 0. Las versiones rápidas
        (especializadas) se prueban siempre contra las generales.
  \item \textbf{Unidades CGS} dentro del núcleo (cm, g, s); conversión a SI solo al mostrar.
  \item \textbf{Determinismo}: mismo montaje, mismo paso de tiempo, mismo resultado. Nada de \texttt{Math.random}
        sin semilla en el núcleo.
  \item \textbf{Pruebas primero}: cada error encontrado se convierte en una prueba nueva antes de corregirlo.
\end{enumerate}

% =================================================================================================
\section{Git: ramas, commits y publicación}

Cada módulo vive en su rama hasta que está cerrado:
\begin{lstlisting}
git checkout main && git pull                 # partir de lo último
git checkout -b modulo-01-cinematica          # rama del módulo
# ... estudiar, implementar ...
node --test "pruebas/*.test.js"               # desde taller/: todo en verde
git add taller && git commit -m "Módulo 1: motor entre rectas e imantado"
git push -u origin modulo-01-cinematica       # copia de seguridad en GitHub
\end{lstlisting}

Al cerrar el módulo:
\begin{lstlisting}
git checkout main && git merge modulo-01-cinematica
git push                                      # GitHub Pages publica /taller/ en 1-2 minutos
\end{lstlisting}

\begin{nota}{Mensajes de commit}
En español y en presente, diciendo qué cambia: «Módulo 2: integrador semiimplícito en el grupo». Un commit por idea:
es más fácil volver atrás y se lee como un diario del aprendizaje.
\end{nota}

% =================================================================================================
\section{Cómo trabajar con Claude en el taller}

\begin{itemize}[leftmargin=*]
  \item \textbf{Contrato del módulo} (paso 5): «genera el contrato del módulo 2». Claude escribe las pruebas y el
        esqueleto a partir de la ficha y de tu código de los módulos anteriores, y valida las pruebas contra una
        implementación de referencia que queda fuera del proyecto, en \texttt{\textasciitilde/Documents/Taller\_referencia}.
  \item \textbf{Revisión}: «revisa mi \texttt{pga.js}». Claude explica cada corrección con la ecuación detrás, sin
        reescribir tu código salvo que lo pidas.
  \item \textbf{Dudas de teoría}: preguntar por la ecuación concreta de la ficha; si la ficha tiene un error, se
        corrige la ficha.
  \item \textbf{La referencia} existe para comparar \emph{después} de intentarlo, no antes.
\end{itemize}

% =================================================================================================
\section{La bitácora}

En \texttt{estudio/bitacora.md}, una entrada por sesión de trabajo:
\begin{lstlisting}
## 2026-11-02 · Módulo 0 · 2 h
- Leído: Doran y Lasenby §2.1-2.4.
- Hecho: ejercicios 0.1 y 0.2. El signo del rotor: e1 gira hacia e2.
- Duda: por qué el regresivo necesita el dual de Poincaré.
- Siguiente: ejercicio 0.3 y productoBlades.
\end{lstlisting}
Sirve para retomar sin perder el hilo y, al final, es el registro de todo el proyecto.

% =================================================================================================
\section{Checklist de cierre de cada módulo}

\begin{itemize}[leftmargin=*]
  \item[{\faSquare[regular]}] Lecturas hechas y ejercicios resueltos.
  \item[{\faSquare[regular]}] Preguntas de control respondidas por escrito.
  \item[{\faSquare[regular]}] Pseudocódigo de la discretización en las soluciones.
  \item[{\faSquare[regular]}] \texttt{node --test} en verde, incluidas las pruebas de módulos anteriores.
  \item[{\faSquare[regular]}] La función pública funciona en computador y celular.
  \item[{\faSquare[regular]}] Rama fusionada con \texttt{main} y subida.
  \item[{\faSquare[regular]}] Entrada de cierre en la bitácora.
\end{itemize}

\end{document}
